\section{Appendix: Experimental Details and Reproducibility}
\label{sec:experimental}
\label{sec:reproducibility}

The \textsc{VOILA} inference procedure is given by the exact utility argmax
in \S\ref{sec:method}.
This appendix provides complete experimental details to ensure reproducibility of all reported results, including feature extraction, calibration procedures, hyperparameters, data splits, cross-validation, and VOI threshold selection.

\subsection{Question Feature Extraction}
\label{sec:features}

All probability estimates used by \textsc{VOILA} are computed from \emph{question-only} features, without access to image pixels or model internal activations.
For each question $q$, we construct a feature vector $\phi(q)$ by concatenating:

\begin{itemize}
    \item \textbf{Lexical TF--IDF features}: We apply a standard TF--IDF vectorizer to the question text after lowercasing and whitespace tokenization.
    The TF--IDF vocabulary is learned \emph{only on the training folds} within each cross-validation split.
    No stopword removal or stemming is applied.
    The resulting TF--IDF dimensionality is fixed within each fold and shared across fidelities.
    
    \item \textbf{Structured question features}: Numeric indicators derived from the question string, including token length, presence of numeric tokens, and binary indicators for common question forms (e.g., yes/no, counting, color, and location queries).
\end{itemize}

TF--IDF and structured features are concatenated and used as input to all fidelity-specific predictors.
Missing feature values are filled with zero, corresponding to the absence of the associated lexical or structural signal.

\subsection{Calibration Model and Hyperparameters}
\label{sec:hyperparams}

For each fidelity level $f$, we train a separate probabilistic correctness predictor using a two-stage calibration pipeline.

\paragraph{Gradient Boosted Regressor (GBR).}
We train a Gradient Boosting Regressor to predict binary correctness labels from $\phi(q)$.
Unless otherwise specified, the default hyperparameters are:
\begin{itemize}
    \item Number of estimators: 100
    \item Learning rate: 0.1
    \item Maximum tree depth: 3
    \item Minimum samples per split: 2
\end{itemize}
All models are trained with a fixed random seed for reproducibility.

\paragraph{Isotonic Calibration.}
The raw GBR outputs are passed through isotonic regression to obtain calibrated correctness probabilities.
Isotonic regression is trained independently for each fidelity, using GBR predictions as input and binary correctness labels as targets.
Out-of-range predictions are handled via clipping.

Within each fold, isotonic calibration is trained exclusively on the training portion of that fold.

\subsection{Cross-Validation Protocol}
\label{sec:cross_validation}

All experiments use \textbf{5-fold cross-validation} over questions.
Specifically, the dataset is partitioned into five disjoint folds at the question level.
For each fold:
\begin{itemize}
    \item Four folds (80\%) are used for training and calibration.
    \item The remaining fold (20\%) is held out for evaluation.
\end{itemize}

This procedure is repeated for all five folds, and all reported results correspond to the \emph{average performance across folds}.
No test fold is ever used for feature extraction, calibration, hyperparameter tuning, or threshold selection.

When standard dataset splits are available (e.g., VQA-v2), we follow the official annotations but perform cross-validation internally to ensure consistent calibration and evaluation across datasets.

\subsection{VOI Policy and Threshold Selection}
\label{sec:voi_policy}

At inference time, \textsc{VOILA} applies an incremental value-of-information (VOI) policy that determines whether to escalate from a lower-cost fidelity $f$ to a higher-cost fidelity $f'$.
Escalation occurs if:
\[
\mathrm{VOI}(f \rightarrow f') = \hat{p}_{f'}(q) - \hat{p}_{f}(q) - \lambda \cdot c(f') > \tau,
\]
where $\hat{p}_f(q)$ denotes the calibrated probability of correctness, $c(f')$ is the acquisition cost, $\lambda$ is a cost scaling parameter, and $\tau$ is a decision threshold.

The parameters $(\lambda, \tau)$ are selected via grid search using \emph{only the training folds} within each cross-validation split.
For each candidate configuration, we optimize a composite cost--accuracy objective that favors dominance over fixed-fidelity baselines across multiple normalized tradeoff metrics.
Selected parameters are then fixed and evaluated on the held-out fold.

\subsection{Hyperparameter Search}
\label{sec:hyperparam_search}

Hyperparameter search jointly explores:
\begin{itemize}
    \item GBR capacity parameters (number of estimators, depth, learning rate),
    \item isotonic regression settings,
    \item VOI decision parameters $(\lambda, \tau)$.
\end{itemize}

Hyperparameters are selected independently within each cross-validation fold using training data only.
Final reported metrics correspond to averages across folds, ensuring robustness to data partitioning.

\subsection{Variance and Reporting}
\label{sec:variance}

Using 5-fold cross-validation allows us to account for variability due to data sampling without relying on repeated random seeds.
We report mean accuracy and average acquisition cost across folds.
Observed performance gaps between \textsc{VOILA} and fixed-fidelity baselines are consistent across folds and substantially larger than fold-to-fold variance.

As our primary goal is to characterize dominance in the cost--accuracy tradeoff rather than marginal accuracy differences, we omit confidence intervals from the main tables for clarity.

Code Availability. We will release the full implementation of VOILA, including preprocessing scripts, calibration code,
and evaluation pipelines, upon acceptance to facilitate reproducibility and further research.



\subsection{Computational Overhead}
\label{sec:computational_overhead}

VOILA introduces minimal computational overhead compared to the costs it optimizes. We measure training and inference times on a single NVIDIA RTX 5000 Ada Generation GPU with 32GB memory using VQA-v2 as a representative benchmark. VLM inference times are measured on 4× V100-SXM2-32GB GPUs.

\paragraph{Training Cost.} 
Training a single configuration of fidelity-specific predictors (GBR + isotonic calibration) requires approximately 21 seconds (0.35 minutes) on VQA-v2 per fidelity. This includes gradient boosting regression training and isotonic calibration for all fidelities. 

During development, hyperparameter search explores combinations of GBR capacity (number of estimators, depth, learning rate), isotonic regression settings, and VOI decision parameters $(\lambda, \tau)$ using 5-fold cross-validation (§\ref{sec:hyperparam_search}). With a typical search space of 20-30 configurations, complete hyperparameter tuning requires 7-10 minutes per dataset. Crucially, this is a one-time offline cost: once trained, the same predictors generalize across query distributions without retraining.

\paragraph{Inference Latency.}
At inference time, VOILA's fidelity selection operates entirely on question features before any visual retrieval or VLM execution. Table~\ref{tab:inference_latency} reports per-query latency averaged over 1000 queries from VQA-v2:

\begin{table}[h]
\centering
\small
\begin{tabular}{lrr}
\toprule
\textbf{Operation} & \textbf{Time (ms)} & \textbf{\% of Total} \\
\midrule
Feature extraction & 0.242 & 53.4\% \\
\quad Basic features (length, type indicators) & 0.025 & 5.5\% \\
\quad TF-IDF transform & 0.217 & 47.9\% \\
GBR forward pass (5 fidelities) & 0.020 & 4.4\% \\
Isotonic calibration (5 fidelities) & 0.030 & 6.6\% \\
VOI decision logic & 0.161 & 35.6\% \\
\midrule
\textbf{Total VOILA overhead} & \textbf{0.453} & \textbf{100\%} \\
\bottomrule
\end{tabular}
\caption{VOILA per-query latency breakdown. Total overhead: 0.453ms. Throughput: 4,132 queries/second on single CPU core.}
\label{tab:inference_latency}
\end{table}

Tables~\ref{tab:vlm_comparison} and~\ref{tab:retrieval_latency} compares VOILA's overhead against VLM inference and retrieval costs. VOILA's 0.45ms overhead is 5,200--140,000× smaller than VLM inference (2.4s--62.8s) and 11--398× smaller than retrieval latency (5--180ms). Feature extraction dominates VOILA's inference cost (53\%), but remains negligible in absolute terms.

\begin{table}[h]
\centering
\small
\begin{tabular}{lrrrr}
\toprule
\textbf{Model} & \textbf{Caption} & \textbf{JPEG Q10} & \textbf{Full-res} & \textbf{Speedup} \\
\midrule
LLaVA-1.5-7B-hf & 2.43 & 5.70 & 5.92 & 5{,}200$\times$ \\
Qwen2.5-VL-7B-Instruct & 2.35 & 5.50 & 5.70 & 5{,}100$\times$ \\
Pixtral-12B-2409 & 3.8 & 7.9 & 8.2 & 8{,}400$\times$ \\
Llama-4-Maverick (17B$\times$12E) & 6.2 & 16.8 & 17.5 & 13{,}700$\times$ \\
Qwen2-VL-72B-Instruct & 4.94 & 25.10 & 25.11 & 10{,}900$\times$ \\
Qwen3-VL-235B-A22B-Instruct & 12.5 & 62.5 & 62.8 & 27{,}700$\times$ \\
\bottomrule
\end{tabular}
\caption{VLM inference latency (seconds) across input fidelities.
Measurements are averaged over 100 runs per fidelity on 4$\times$ V100-SXM2-32GB GPUs.
Speedup is computed relative to \textsc{VOILA} overhead.}
\label{tab:vlm_comparison}
\end{table}

\begin{table}[h]
\centering
\small
\begin{tabular}{lrr}
\toprule
\textbf{Operation} & \textbf{Latency} & \textbf{Speedup} \\
\midrule
Edge-cloud (caption, 2KB) & 5 ms & 11$\times$ \\
Edge-cloud (JPEG Q10, 45KB) & 65 ms & 144$\times$ \\
Edge-cloud (full-res, 650KB) & 180 ms & 398$\times$ \\
\midrule
\textbf{VOILA overhead} & \textbf{0.453 ms} & \textbf{---} \\
\bottomrule
\end{tabular}
\caption{Retrieval latency and \textsc{VOILA} overhead.
Retrieval measured in an edge--cloud deployment with 4G LTE (20\,Mbps uplink) to AWS S3.
Speedup denotes the ratio of full-resolution retrieval latency to \textsc{VOILA} overhead.}
\label{tab:retrieval_latency}
\end{table}



Critically, VOILA's overhead is constant regardless of selected fidelity or VLM size: the same 0.45ms cost applies whether routing to caption (2KB, 2.4s inference) or full-resolution (650KB, 62.8s inference). This enables cost-effective adaptive routing even for the largest models.

\paragraph{Memory Footprint.}
The complete system requires only 0.04MB for all learned components (5 fidelity-specific GBR models + isotonic calibrators). This minimal footprint enables deployment on resource-constrained edge devices, smartphones, and IoT systems. No GPU memory is required at inference time, as all computation operates on CPU using lightweight tree ensemble inference and isotonic lookup tables.

%\paragraph{Scalability.}
%VOILA's computational cost scales as $O(K \cdot d \cdot \log T)$ per query, where $K$ is the number of fidelities (typically 3--7), $d$ is feature dimensionality (2,000--5,000 for TF-IDF), and $T$ is the number of trees in the GBR ensemble (typically 100). This remains negligible even as $K$ and $d$ grow, and is independent of VLM size, image resolution, and context length. The system processes over 4,000 queries per second on a single CPU core, orders of magnitude faster than VLM inference or network retrieval.